\chapter{Fazit}

Es wurde anhand einer Fast-Slow-Koinzidenzmessung mit Nuclear Instument Modulen (NIMs) die Lebensdauer des $(\frac{5}{2}^+)$-Zustands von \caesium untersucht. Dazu wurden mithilfe von \na- und \ba-Quellen die einzelnen Teile des Schaltkreises für die Messung (siehe Abbildung \ref{fig:tikz/fast_slow_koinzidenz}) für beide Seiten des Aufbaus (mit jeweils einem der beiden NaI(Tl)-Szintillationsdetektoren) justiert und auf die Start- und Stopp-Photonen für die Koinzidenz mit den Energien $E_\text{Start}=\SI{356.0}{\kilo\eV}$ und $E_\text{Stopp}=\SI{81.0}{\kilo\eV}$ eingestellt.

Als erstes wurden die Hoch- und Niedervolt-Versorgungsspannungen des Aufbaus kontrolliert. Hierbei fiel auf, dass die $\SI{-6}{\V}$ und $\SI{6}{\V}$ Anschlüsse der oberen Reihe nicht wie erwartet funktionierten, sondern vermutlich einen defekten Gleichrichter aufwiesen. Da diese für die verwendeten Module aber nicht gebraucht wurden, hatte dies keine Konsequenzen auf die weiteren Messungen.
Anschließend wurden die Verstärkungen und Verzögerungen im Slow-Kreis angepasst, sowie mithilfe einer Aufnahme des gesamten Energiespektrums der \na-Quelle (mit ganz offenen SCA-Fenstern) die Fenster der SCAs auf die $\SI{511}{\kilo\eV}$-Annihilationslinie eingestellt. Vor Einstellung des Fensters auf die \na-Linie wurde mit der \ba-Quelle zudem auch noch das gesamte mögliche Energiespektrum dieser Quelle aufgenommen, indem die SCA-Fenster komplett geöffnet wurden. Anhand der in diesen beiden Spektren sichtbaren Photopeaks wurden diesen bekannte Energien von $\gamma$-Photonen aus den beiden Spektren (vgl. Abbildungen \ref{fig:zerfallbarium} und \ref{fig:zerfallnatrium}) zugeordnet. An diese Zuordnung wurde eine Geradengleichung angepasst, um den Zusammenhang zwischen MCA-Kanälen (Bins) und Energien herzustellen. Die Anpassungsgüte war hier sehr hoch, die Kalibration lief sehr gut.

Es wurde mithilfe der Verzögerung an den GDGs der Slow-Zweige der beiden Seiten mit der \na-Quelle eine zeitliche Überlappung der Signale hergestellt, um die Koinzidenz sicherzustellen.

Anschließend wurden die Verzögerungen und die Schwellen der CFDs in den beiden Fast-Kreisen angepasst, um die Signale der beiden Photonen zugehörig zum untersuchten $\frac{5}{2}^+$-\caesium-Zustand in der richtigen Reihenfolge am TAC ankommen zu lassen. Es wurde überprüft, dass die Signale am Ende des Slow-Kreises (nach der Koinzidenzeinheit) und am Ende des Fast-Kreises (Ausgang des TAC) grundsätzlich koinzident waren, sie sich also überlappten. Dies konnte festgestellt werden.

Zusätzlich wurde eine Zeitkalibration für den TAC durchgeführt, dabei wurde mit der \na-Quelle, bei dem die beiden $\SI{511}{\kilo\eV}$-Photonen immer zeitgleich sind, und variabler Verzögerung mit dem ns-Delay-Bauteil in einem der beiden Fast-Kreise das Energiespektrum mit auf diese Energie eingestellte SCA-Fenster aufgenommen (Promptkurve). Die Verzögerung wurde in konstanten Schritten erhöht, es wurden fünf Promptkurven aufgenommen. Die eingestellten Verzögerungswerte wurden gegen die Positionen der Schwerpunkte von an diese angepassten Gaußfunktionen aufgetragen und an diesen Zusammenhang eine Geradengleichung angepasst. Diese Geradengleichung ergab so einen Zusammehnag zwischen MCA-Kanälen (Bins) und Zeitverzögerung der beiden Photonenpulse. Diese Anpassung zeigte ein sehr gutes Maß der Anpassungsgüte und konnte so zur Umrechnung von Bins zu (Verzögerungs-)Zeit genutzt werden.

Es wurden dann noch die Fenster der SCAs der beiden Seiten für die \ba-Quelle auf die gesuchten Linien $\SI{81.0}{\kilo\eV}$ und $\SI{356.0}{\kilo\eV}$ eingestellt und dann mit dieser Quelle eine Lebensdauerkurve aufgenommen.

Durch die Messung der Zeitdifferenzen der Signale der Start- und Stopp-Photonen konnte dann die Lebensdauerkurve (Abbildung \ref{fig:lifetime_plot}) des zu untersuchenden \caesium Zustands aufgenommen werden. Aus der Form der Kurve wurden dann die mittlere Lebensdauer und die Halbwertszeit als \chapter{Lebensdauermessung}

\section{Zeitkalibration}
Um die Zeitkalibration und informationen über die Zeitauflösung gewinnen zu können, werden an die Promptkurven wieder Gaußfunktionen angepasst. Die Ergebnisse dieser Anpassungen sind in \fit{prompt} und \tab{prompt} zu sehen. Die einzelnen Gaußanpassungen weisen hierbei sehr gute Fitgüten auf.

\jafps{prompt}{Promptkurve mit angepassten Gaußfunktionen.}{0.75}
\jaftl{|c|c|c|c|c|}{../data/fits/prompt.table}{}{prompt}

Die so gewonnenen Positionen werden nun Aufgetragen, gegen die eingestellten Zeiten. Hierbei ist nur der Zeitunterschied von $\SI{16}{\nano\second}$ zwischen den einzelnen Kurven relevant, also Startwert wurden allerdings $\SI{8}{\nano\second}$ gewählt, da so stets die Einstellungen des ns-Delay genau die aufgetragene Zeit sind. An diese Auftragung wird schließlich eine Grade angepasst. Dies ist in \fig{time_cal} abgebildet. Die Parameter der Gradenanpassung sind dann \input{../data/fits/time_cal_fit.txt}. Die Anpassungsgüte ist hierbei sehr hoch, was für eine ausgezeichnete Linearität der Zeitmessung spricht.

Aus den Standardabweichungen der Gaußkurven in \tab{prompt} kann nun die Halbwertsbreite (Full-Width-Half-Maximum (FWHM)) als Maß der Breite der Kurven genutzt werden. Dabei ist $\text{FWHM} = \sigma \cdot \sqrt{8 \ln{2}}$ \cite[8]{ex_techniques}. $\text{FWHM}_\text{ns} = \text{FWHM} \cdot a$, wobei die Fehlerfortpflanzung \todo{adsflkjasdasdfa} ist.

Mittels der soeben bestimmten Gradengleichung kann diese Breite auf eine Zeit umgerechnet werden. Somit erhählt man ein Maß für die Zeitauflösung des Gesamtsystems. Die Berechnungen für die einzelnen Kurven sind in \tab{time_resolution.table} aufgeführt. Im Mittel erhält man dann eine Halbwertsbreite von \input{../data/fits/fwhm_time_res.txt}. Diese kann als Maß für die Zeitauflösung genutzt werden. Diese Zeitauflösung ist nur in der gleichen Größenordnung wie die vermutete Lebensdauer, welche gemessen werden soll. Würde man sich bei der Lebensdauermessung also nur auf die reine Unterscheidbarkeit zweier Linien stützen, so wäre die Messung mitunter nicht durchführbar. Wie im folgenden jedoch erläutert wird bezieht die Finale Messung diese Punktspreizfunktion (PSF) mit ein.


\jaft{|c|c|c|}{time_resolution.table}{\todo{caption}}


\jafps{time_cal}{Zeitkalibration}{0.9}


\section{Lebnsdauer eines Angeregten Cäsiumzustandes}

Um die Lebensdauer des $\frac{5}{2}^+$ \caesium Zustandes zu bestimmen soll an die Lebensdauerkurve eine Zerfallskurve angepasst werden. Da jedoch, eine einzelne Zeit durch die PSF Gaußförmig verschmiert wird, und sich dieser Prozess mathematisch als Faltung beschreiben lässt muss insgesamt die Faltung einer Zerfallskurve und einer Gaußkurve angepasst werden. Die insgesamte Funktion ist dann \cite{525va}

\begin{align}
	1
\end{align}


 bestimmt. Diese Halbwertszeit weicht vom Literaturwert um etwa eine halbe ns ab, ist damit aber in der korrekten Größenordnng. Bei Beachtung der systematischen Unsicherheiten des Systems wird aufgrund der endlichen Energieauflösung des Aufbaus klar, dass diese die Energie des Start-Photons nicht gut genug auflösen kann, um eine Unterscheidung zum ebenfalls als Zerfallskanal auftretenden Photon mit $E_\text{alt}=\SI{383.8}{\kilo\eV}$ zu ermöglichen. Dies erklärt die Überschätzung der mittleren Halbwertszeit und damit der Lebensdauer. Mit einer niedrigeren oberen SCA-Schwelle für das Startphoton könnte dieses Problem verbessert werden, dafür müsste die Messdauer erhöht werden.
